\section{Proof dependencies and the complete remaining source}
\label{app:v104-dependencies}
The original event and target are \eqref{eq:v76-event} and
\eqref{eq:v76-target}. The inherited inputs include the complete
positive paired measure, its fixed-width concentration, actual uncut
optical volume, weighted variation, regular critical trace and the
original scalar schedule. Their historical proofs remain in the
archive; no new spectral or pressure theorem is assumed.

Endpoint inclusion and integrated positive-trace balance give
Theorem~\ref{thm:v104-block} at finite collision count. The critical
incoming normalization and the width of one actual comparison interval
give Lemma~\ref{lem:v104-arm-balance}. If the outgoing trace is too
small, Theorem~\ref{thm:v104-paid-arms} pays its entire shortage by
the original uncut volume in one common roof window. The resulting
height estimate needs no outgoing trace-fraction lower bound.
Corollary~\ref{cor:v104-whole-arm-recovery} then removes the entire
selected critical-arm contribution from the old three-clock failure.
Physical overlap is counted on the disjoint angular sheet space,
not by counting words or strata.

The finite-order condition is local to an actual arm. Its comparison
levels must be reached before any true boundary, and its geometric
ratio must meet the stated uniform or outer-budget bound. No uniform
order, reach or complete physical coverage is inferred from
finite-count definability. The stationary model checks the distinction
from the old clocks; it is not an additional Lorentz realization.

The complete remaining source is positive and lies outside the
recovered physical families. Vanishing of that source in the original
essential-height norm is still required for the organizing theorem.
The complete preceding article and all earlier mathematical cores
are reproduced in the supplement, with the two revision-88 variants
separately identified and preserved.
